\documentclass[12pt,a4paper]{report}
\usepackage{fontspec}
\IfFontExistsTF{TeX Gyre Termes}{\setmainfont{TeX Gyre Termes}}{\setmainfont{Liberation Serif}}
\IfFontExistsTF{TeX Gyre Heros}{\setsansfont{TeX Gyre Heros}}{\setsansfont{Liberation Sans}}
\setmonofont[Scale=0.84]{DejaVu Sans Mono}
\usepackage[vietnamese]{babel}
\usepackage[a4paper,margin=1in]{geometry}
\usepackage{amsmath,amssymb,mathtools,bm}
\usepackage{booktabs,tabularx,array,longtable}
\usepackage{graphicx,float}
\usepackage[hidelinks]{hyperref}
\usepackage{fancyhdr,setspace,titlesec,caption,microtype,needspace}
\usepackage{enumitem}
\setstretch{1.15}
\setlength{\parskip}{0.3em}
\setlength{\parindent}{0pt}
\setlength{\headheight}{16pt}
\setcounter{tocdepth}{1}
\setcounter{secnumdepth}{2}
\captionsetup{font=small,labelfont=bf}
\numberwithin{equation}{chapter}
\titleformat{\chapter}[display]{\normalfont\bfseries}{\filleft\Large Chương \thechapter}{0.3em}{\filright\huge}[\vspace{0.4em}\titlerule]
\titlespacing*{\chapter}{0pt}{-10pt}{20pt}
\titleformat{\section}{\large\bfseries}{\thesection}{0.6em}{}
\titleformat{\subsection}{\normalsize\bfseries}{\thesubsection}{0.6em}{}
\pagestyle{fancy}
\fancyhf{}
\lhead{\small CO5085 · Bài tập E1–E3}
\rhead{\small Huỳnh Thành Đạt · 2570161}
\cfoot{\thepage}
\fancypagestyle{plain}{\fancyhf{}\lhead{\small CO5085 · Bài tập E1–E3}\rhead{\small Huỳnh Thành Đạt · 2570161}\cfoot{\thepage}}
\newcommand{\R}{\mathbb{R}}
\newcommand{\norm}[1]{\left\lVert #1 \right\rVert}
\newcommand{\code}[1]{\texttt{\detokenize{#1}}}
\newcommand{\resultfigure}[3]{\IfFileExists{generated/figures/#1.png}{\begin{figure}[htbp]\centering\includegraphics[width=#3\textwidth]{generated/figures/#1.png}\caption{#2}\end{figure}}{}}
\newcommand{\tablenote}{\par\smallskip{\small Val/Test: accuracy (\%). Tham số đếm từ kiến trúc; $\pm$: SD qua ba seed. Các lượt chạy được liệt kê ở Chương 6; cách đo thời gian ở Chương 1.}\par}
\hypersetup{pdftitle={CO5085 - Bao cao E1-E3},pdfauthor={Huynh Thanh Dat - 2570161}}
\emergencystretch=2em
\begin{document}
\hypersetup{pageanchor=false}
\begin{titlepage}
\centering
{\large ĐẠI HỌC QUỐC GIA THÀNH PHỐ HỒ CHÍ MINH\par}
{\large\bfseries TRƯỜNG ĐẠI HỌC BÁCH KHOA\par}
{\large KHOA KHOA HỌC VÀ KỸ THUẬT MÁY TÍNH\par}
\vspace{1.2cm}
\includegraphics[width=3cm]{assets/hcmut-logo.png}
\vspace{0.8cm}

{\large\bfseries HỌC SÂU VÀ ỨNG DỤNG\par}
{\large\bfseries TRONG THỊ GIÁC MÁY TÍNH\par}
{\large Mã môn học: CO5085\par}
\vspace{0.9cm}
\rule{\textwidth}{0.6pt}\par
\vspace{0.35cm}
{\Huge\bfseries BÁO CÁO BÀI TẬP E1–E3\par}
\vspace{0.45cm}
{\Large Phân loại ảnh với CNN, attention\par}
{\Large và mạng hồi quy\par}
\vspace{0.35cm}
\rule{\textwidth}{0.6pt}\par
\vspace{0.7cm}
{\large Học kỳ 261 · Năm học 2026–2027\par}
\vspace{0.7cm}
\begin{tabular}{@{}ll@{}}
Giảng viên: & Lê Thành Sách\\
& Nguyễn Quốc Minh\\[0.5em]
Sinh viên: & Huỳnh Thành Đạt\\
Mã số sinh viên: & 2570161\\
Hình thức: & Cá nhân
\end{tabular}
\vfill
{\normalsize TP. Hồ Chí Minh, tháng 10 năm 2026\par}
\end{titlepage}
\hypersetup{pageanchor=true}
\pagenumbering{roman}
{\setstretch{1.0}\small\tableofcontents}
\clearpage
\pagenumbering{arabic}

\chapter{Dữ liệu và quy trình chung}
\section{Phạm vi bài làm}
Báo cáo so sánh các mô hình phân loại ảnh trên cùng tập Fashion-MNIST. E1 gồm softmax, MLP và CNN. E2 trình bày hiện thực multi-head self-attention (MSA), đối chiếu với thư viện và khảo sát ba cách tạo token. E3 biểu diễn ảnh thành chuỗi để phân loại bằng LSTM/GRU, với MLP và CNN của E1 làm đối chứng. Chỉ số bước của mạng hồi quy tương ứng với thứ tự quét các vùng trên ảnh.

Phần mở rộng khảo sát bốn thay đổi trên LSTM: khởi tạo ma trận hồi quy trực giao, regularization trực giao mềm, khởi tạo bias cổng theo độ dài chuỗi và mean pooling. Mỗi thành phần được đánh giá riêng, sau đó kết hợp và thực hiện ablation. Tất cả mô hình trong E1–E3 đều được huấn luyện từ đầu.

Thực nghiệm gồm 20 cấu hình, mỗi cấu hình chạy với ba seed 42, 43 và 44, tương ứng 60 lượt huấn luyện. Các mô hình dùng chung cách chia dữ liệu và quy tắc chọn checkpoint để so sánh kết quả.

\section{Tập dữ liệu và cách chia}
Fashion-MNIST gồm ảnh thời trang mức xám $28\times28$, 10 lớp, với 60.000 ảnh train và 10.000 ảnh test chính thức \cite{fashion}. Kích thước ảnh nhỏ phù hợp để khảo sát cả CNN, attention và mạng hồi quy trên cùng đầu vào.

Từ 60.000 ảnh train, tách phân tầng 10\% làm validation với split seed 2026. Như vậy có 54.000 ảnh train, 6.000 validation và 10.000 test. Mỗi lớp lần lượt có 5.400, 600 và 1.000 ảnh theo phân bố chuẩn của bộ dữ liệu. Split seed không đổi giữa E1–E3 và không trùng vai trò với seed khởi tạo mô hình.

Ảnh được chuyển từ uint8 sang $[0,1]$, sau đó chuẩn hóa $x'=(x-\mu)/\sigma$. Hai thống kê $\mu,\sigma$ được tính trên 54.000 ảnh train và dùng lại cho validation/test. Các thí nghiệm không dùng augmentation để giữ đầu vào thống nhất khi so sánh cách biểu diễn ảnh.

\code{src/common/data.py} lưu chỉ số từng split, checksum dữ liệu và fingerprint của cả split lẫn normalization. Loader ưu tiên HTTPS từ repo chính thức của dataset và giữ kiểm tra checksum của torchvision. Dữ liệu test có thể được tải cùng kho ảnh, nhưng không được đọc vào vòng chọn checkpoint hay tính normalization.

\section{Huấn luyện và chỉ số đánh giá}
Mỗi mô hình xuất logits $z\in\R^{10}$. Với nhãn đúng $y$, cross-entropy là
\begin{equation}
\ell(z,y)=-z_y+\log\sum_{k=1}^{10}\exp(z_k),
\qquad \frac{\partial\ell}{\partial z_k}=p_k-\mathbf{1}[k=y].
\end{equation}
Vì \code{cross_entropy} đã thực hiện log-softmax ổn định số, code không chèn softmax trước loss. Softmax chỉ được dùng khi cần xác suất để trình bày dự đoán.

\begin{table}[htbp]\centering
\begin{tabularx}{\textwidth}{@{}lX@{}}\toprule
Thiết lập & Giá trị định trước\\\midrule
Optimizer & AdamW, weight decay $10^{-4}$ trên các tham số trainable\\
Ngân sách & 30 epoch, batch 256, không early stopping\\
Learning rate & E1/E3: $10^{-3}$; E2: $3\times10^{-4}$\\
Scheduler & Cosine decay về 1\% learning rate ban đầu\\
Gradient clipping & Global norm tối đa 1, sau backward và trước step\\
Độ chính xác số & FP32; tắt AMP, TF32 và compile\\
Seed mô hình & 42, 43, 44; cùng split seed 2026\\
Chọn checkpoint & Validation accuracy lớn nhất; hòa thì CE nhỏ nhất\\\bottomrule
\end{tabularx}
\caption{Thiết lập huấn luyện dùng chung cho E1–E3.}
\end{table}

Vòng lặp huấn luyện gồm zero gradient, forward, tính loss, backward, clipping và optimizer step. Checkpoint được chọn trên validation; test được đánh giá sau khi hoàn tất các cấu hình đã định. Chương trình kiểm tra tính hữu hạn của loss và gradient tại mỗi bước cập nhật.

Accuracy là chỉ số chính; macro-F1 bổ sung đánh giá theo lớp. Confusion matrix và các dự đoán sai có confidence cao được dùng để phân tích lỗi. Chi phí mô hình được mô tả qua số tham số và thời gian train/epoch. Thời gian được đồng bộ CUDA, gồm lấy batch và cập nhật trọng số; median bỏ epoch đầu để giảm ảnh hưởng khởi tạo. Tốc độ inference được đo riêng, bỏ hai batch đầu và không tính thời gian truyền dữ liệu sang GPU.

Bảng kết quả trình bày trung bình và độ lệch chuẩn mẫu qua ba seed. Các đường cong, confusion matrix và ảnh minh họa dùng seed 42. Ngân sách 30 epoch được giữ cố định cho mọi cấu hình; do đó kết quả phản ánh hiệu quả trong ngân sách này.

\chapter{E1: Softmax, MLP và CNN}
\section{Ba bộ phân loại}
Softmax nhận vector $x\in\R^{784}$ và tính logits $z=Wx+b$; ranh giới giữa hai lớp là siêu phẳng trong không gian pixel. MLP bổ sung hai tầng phi tuyến trên cùng vector flatten. CNN dùng kernel chia sẻ theo vị trí để học đặc trưng cục bộ. Ba kiến trúc được giữ nhỏ và mô tả đầy đủ trong bảng dưới đây.

\begin{table}[htbp]\centering
\begin{tabularx}{\textwidth}{@{}lX@{}}\toprule
Mô hình & Các lớp theo thứ tự\\\midrule
Softmax & Flatten $784\to10$\\
MLP & Flatten; $784\to256$; ReLU; dropout 0,2; $256\to128$; ReLU; dropout 0,2; $128\to10$\\
CNN & Conv $1\to32$, kernel $3$, padding $1$; BN; ReLU; pool $2$; Conv $32\to64$ tương tự; pool $2$; flatten; $3136\to128$; ReLU; dropout 0,2; $128\to10$\\\bottomrule
\end{tabularx}
\caption{Kiến trúc E1. Conv không dùng bias vì sau đó có batch normalization.}
\end{table}

Hai phép pooling đưa ảnh $28\to14\to7$, vì vậy head CNN nhận $64\times7\times7=3136$ giá trị. Với một lớp Linear $d_{in}\to d_{out}$ có bias, số tham số là $d_{in}d_{out}+d_{out}$. Ví dụ softmax có $784\times10+10=7.850$ tham số. Các số còn lại được đếm trực tiếp bằng \code{numel()}, gồm tham số affine của BN nhưng không gồm running statistics.

\section{Kết quả và phân tích lỗi}
\begin{table}[htbp]\centering\small
\begin{tabularx}{\textwidth}{@{}Xrrrr@{}}
\toprule
Mô hình & Tham số & Val (\%) & Test (\%) & Giây/epoch \\
\midrule
Softmax & 7,850 & 85.98 & 84.51 $\pm$ 0.15 & 1.51 \\
MLP & 235,146 & 91.02 & 89.88 $\pm$ 0.14 & 1.67 \\
CNN & 421,738 & 93.03 & 92.69 $\pm$ 0.15 & 1.97 \\
\bottomrule
\end{tabularx}

\caption{Kết quả E1, tổng hợp theo các seed được công bố.}
\end{table}
\tablenote

Train/validation CE được đối chiếu với accuracy để phân tích quá trình học. Confusion matrix và ảnh dự đoán sai bên dưới lấy từ CNN ở seed 42.
\IfFileExists{generated/e1_analysis.tex}{Các kết quả được đo với ba seed 42, 43 và 44 trên NVIDIA GeForce RTX 5090; PyTorch 2.10.0+cu130 và CUDA runtime 13.0. Chênh lệch giữa hai cấu hình được tính theo cặp cùng seed, rồi lấy trung bình và SD.

Softmax: test accuracy $(84.51 \pm 0.15)\%$, macro-F1 0.8436; CE cuối epoch train/validation trung bình 0.3663/0.4095.

MLP: test accuracy $(89.88 \pm 0.14)\%$, macro-F1 0.8985; CE cuối epoch train/validation trung bình 0.1371/0.2832.

CNN: test accuracy $(92.69 \pm 0.15)\%$, macro-F1 0.9267; CE cuối epoch train/validation trung bình 0.0174/0.2969.

MLP trừ Softmax: +5.37 $\pm$ 0.29 điểm phần trăm.

CNN trừ MLP: +2.81 $\pm$ 0.01 điểm phần trăm.

MLP thêm biểu diễn phi tuyến; CNN giữ cấu trúc hai chiều, chia sẻ kernel và pooling. Xếp hạng đã đo phù hợp với lợi ích của các biểu diễn này trên ảnh nhỏ. Khác biệt về số tham số và regularization cũng góp phần vào kết quả của từng kiến trúc.

CNN có khoảng cách train/validation CE cuối kỳ lớn dù test accuracy cao nhất E1. Đường CE và error gallery gợi ý overfit hoặc dự đoán sai quá tự tin; calibration chưa được khảo sát riêng trong bài này. Test dùng checkpoint được validation chọn, còn các CE cuối kỳ ở trên mô tả epoch 30.

Ở CNN seed 42, lỗi nhiều nhất theo một chiều là Shirt thành T-shirt/top: 90 ảnh.
}{}

\resultfigure{e1_comparison_ce}{Validation CE của ba mô hình E1 ở seed 42.}{0.88}
\resultfigure{e1_comparison_accuracy}{Validation accuracy của E1; seed dùng cho hình ghi trên biểu đồ.}{0.88}
\resultfigure{e1_cnn_confusion}{Confusion matrix test của CNN ở seed đại diện.}{0.8}
\resultfigure{e1_cnn_mistakes}{Một số dự đoán sai có confidence cao của CNN.}{0.95}

\chapter{E2: MSA và cách lấy token}
\section{Hiện thực attention}
Với tensor token $X\in\R^{B\times N\times d}$, mỗi head học các phép chiếu query, key và value. Bỏ chỉ số batch để gọn, attention của một head là \cite{attention}
\begin{equation}
Q=XW_Q,\quad K=XW_K,\quad V=XW_V,\qquad
A=\operatorname{softmax}_{\text{hàng}}\left(\frac{QK^\top}{\sqrt{d_h}}\right),\quad O=AV.
\end{equation}
Các bias được lược bỏ trong công thức cho gọn, nhưng có trong hiện thực. Ở đây $d=64$, có bốn head nên $d_h=16$. Ba phép Linear được gộp thành QKV, reshape thành các head, nhân ma trận, softmax theo chiều key, áp dụng dropout, ghép head và chiếu output. Bản manual không gọi \code{nn.MultiheadAttention}, \code{TransformerEncoder} hoặc fused attention để thay phần tính này.

Bản đối chiếu dùng \code{nn.MultiheadAttention(batch_first=True)}. Chương trình tạo toàn bộ model manual rồi copy đúng QKV và output projection sang module PyTorch. Các phần tokenizer, positional embedding, FFN, LayerNorm và head cũng giữ nguyên tensor khởi tạo của cùng seed. Chỉ reset seed trước hai constructor khác nhau là chưa đủ bảo đảm điều này.

Tính tương đương số học được kiểm tra bằng tensor nhỏ và dropout bằng 0, qua output, gradient đầu vào và gradient của toàn bộ tham số. Khi huấn luyện, dropout và kernel có thể làm hai phiên bản đi theo quỹ đạo khác nhau dù cùng khởi tạo.

\section{Ba cách tokenization}
\begin{table}[htbp]\centering
\begin{tabularx}{\textwidth}{@{}lrrX@{}}\toprule
Cách chia & $N$ & Đầu vào/token & Phép tạo embedding\\\midrule
Hàng & 28 & 28 & Linear $28\to64$\\
Patch $4\times4$ & 49 & 16 & Unfold không chồng lấp, Linear $16\to64$\\
CNN-stem & 49 & 64 & Hai Conv stride 2: $1\to32\to64$, BN/ReLU; flatten spatial\\\bottomrule
\end{tabularx}
\caption{Tokenization của E2 trên ảnh $1\times28\times28$.}
\end{table}

Token hàng gom các pixel cùng một hàng. Token patch quét lưới từ trái sang phải, trên xuống dưới. CNN-stem rút feature map $64\times7\times7$; mỗi vị trí là một token đã qua xử lý cục bộ. Stem được huấn luyện từ đầu và bổ sung cả phép tính lẫn tham số so với hai cách tạo token trực tiếp.

Mỗi cấu hình thêm learned positional embedding, qua hai encoder block dạng pre-LN và kết thúc bằng LayerNorm, mean pooling và Linear $64\to10$. Không có CLS token. Mỗi block tính
\begin{align}
U &= X+\operatorname{Dropout}(\operatorname{MSA}(\operatorname{LN}(X))),\\
Y &= U+\operatorname{FFN}(\operatorname{LN}(U)),
\end{align}
trong đó FFN là $64\to128\to64$ với GELU và dropout 0,1. Attention không dùng causal mask vì toàn bộ ảnh đã có tại thời điểm phân loại.

\section{So sánh attention và tokenization}
\begin{table}[htbp]\centering\small
\begin{tabularx}{\textwidth}{@{}Xrrrr@{}}
\toprule
Mô hình & Tham số & Val (\%) & Test (\%) & Giây/epoch \\
\midrule
patches / manual & 71,946 & 87.03 & 85.93 $\pm$ 0.60 & 4.52 \\
patches / torch & 71,946 & 87.09 & 86.02 $\pm$ 0.37 & 4.46 \\
rows / manual & 71,370 & 88.61 & 87.92 $\pm$ 0.03 & 4.93 \\
rows / torch & 71,370 & 88.61 & 87.86 $\pm$ 0.02 & 4.60 \\
cnn / manual & 89,770 & 89.26 & 88.66 $\pm$ 0.40 & 5.10 \\
cnn / torch & 89,770 & 89.40 & 88.63 $\pm$ 0.35 & 4.76 \\
\bottomrule
\end{tabularx}

\caption{Sáu cấu hình E2. Manual và PyTorch của cùng tokenizer có cùng số tham số.}
\end{table}
\tablenote

Manual và thư viện được so sánh ở từng tokenizer; ba tokenizer được so sánh khi giữ attention backend cố định. Chi phí phần score của attention tỉ lệ $N^2d$; số cặp token của patch gấp $49^2/28^2\approx3,06$ lần token hàng. Thời gian toàn mô hình còn gồm projection, FFN, stem và chi phí thực thi kernel.

Các kết quả dưới đây được đo trên cùng split, với cùng thiết lập giữa hai backend của mỗi tokenizer.
\IfFileExists{generated/e2_analysis.tex}{Tokenizer rows: PyTorch trừ manual có chênh lệch test -0.06 $\pm$ 0.04 điểm phần trăm; thời gian train/epoch trung bình của median từng seed lần lượt 4.60s và 4.93s.

Tokenizer patches: PyTorch trừ manual có chênh lệch test +0.10 $\pm$ 0.23 điểm phần trăm; thời gian train/epoch trung bình của median từng seed lần lượt 4.46s và 4.52s.

Tokenizer cnn: PyTorch trừ manual có chênh lệch test -0.02 $\pm$ 0.10 điểm phần trăm; thời gian train/epoch trung bình của median từng seed lần lượt 4.76s và 5.10s.

Trong sáu cấu hình, cnn / manual có test accuracy trung bình cao nhất, đạt 88.66\%.

Giữ backend manual, CNN-stem trừ rows / manual: +0.74 $\pm$ 0.42 điểm phần trăm.

Giữ backend manual, CNN-stem trừ patches / manual: +2.73 $\pm$ 0.31 điểm phần trăm.

Cấu hình E2 tốt nhất so với CNN E1 có chênh lệch -4.03 $\pm$ 0.46 điểm phần trăm. CNN-stem cung cấp đặc trưng cục bộ trước attention; các kết quả gợi ý rằng prior không gian hữu ích trong ngân sách này. Tokenizer đồng thời làm thay đổi số token, chiều đầu vào và số tham số.
}{}

\resultfigure{e2_comparison_ce}{Validation CE của sáu cấu hình E2 ở seed 42.}{0.9}
\resultfigure{e2_comparison_accuracy}{Validation accuracy khi thay attention backend và tokenizer.}{0.9}

\chapter{E3: LSTM/GRU trên chuỗi ảnh}
\section{Biểu diễn ảnh và kiến trúc}
Ảnh được biến thành chuỗi theo hai cách đã dùng trong E2, nhưng không qua embedding học được: hàng cho $T=28,D=28$; patch $4\times4$ cho $T=49,D=16$. Cả LSTM và GRU đều có một tầng, một chiều, hidden size $H=128$, trạng thái đầu bằng zero. Baseline lấy hidden state cuối, dropout 0,1 rồi Linear $128\to10$.

Quét hàng giữ các quan hệ ngang trong một token và để mạng hồi quy xử lý quan hệ giữa các hàng. Quét patch tạo nhiều bước hơn, mỗi bước mô tả một vùng nhỏ hơn. Vì vậy, so sánh này thay đổi đồng thời độ dài chuỗi và kích thước đầu vào tại mỗi bước.

\section{Các phương trình được sử dụng}
Với LSTM, đặt $a_t=W_xx_t+W_hh_{t-1}+b$ rồi tách thành bốn phần \cite{pytorch-lstm}:
\begin{align}
i_t&=\sigma(a_t^i),& f_t&=\sigma(a_t^f),& g_t&=\tanh(a_t^g),& o_t&=\sigma(a_t^o),\\
c_t&=f_t\odot c_{t-1}+i_t\odot g_t,&&
h_t=o_t\odot\tanh(c_t).
\end{align}
$c_t$ có đường cập nhật cộng, còn $h_t$ là trạng thái dùng cho output và hồi tiếp. Cổng quên lớn có thể giữ thông tin, nhưng kiến trúc này không tự bảo đảm mọi gradient đều được bảo toàn.

GRU dùng reset gate $r_t$, update gate $z_t$ và trạng thái ứng viên $n_t$. Diagnostic khớp đúng cách PyTorch đặt reset gate sau phép chiếu hidden \cite{pytorch-gru}:
\begin{align}
r_t &=\sigma(W_{ir}x_t+b_{ir}+W_{hr}h_{t-1}+b_{hr}),\\
z_t &=\sigma(W_{iz}x_t+b_{iz}+W_{hz}h_{t-1}+b_{hz}),\\
n_t &=\tanh\!\left(W_{in}x_t+b_{in}+r_t\odot(W_{hn}h_{t-1}+b_{hn})\right),\\
h_t &=(1-z_t)\odot n_t+z_t\odot h_{t-1}.
\end{align}
Trong công thức của PyTorch, reset gate tác động sau phép chiếu hidden. Đây là điểm khác với cách viết $W_{hn}(r_t\odot h_{t-1})$. GRU sử dụng một trạng thái hidden, còn LSTM tách hidden state và cell state.

PyTorch lưu hai vector bias cho mỗi tầng. Vì vậy, số tham số phần recurrent là $4H(D+H+2)$ với LSTM và $3H(D+H+2)$ với GRU; cộng thêm $10H+10$ cho head. Giữ cùng $H$ giúp so sánh rõ nhưng không làm hai mô hình có cùng capacity.

\section{Kết quả và quan hệ không gian}
\begin{table}[htbp]\centering\small
\begin{tabularx}{\textwidth}{@{}Xrrrr@{}}
\toprule
Mô hình & Tham số & Val (\%) & Test (\%) & Giây/epoch \\
\midrule
MLP & 235,146 & 91.02 & 89.88 $\pm$ 0.14 & 1.67 \\
CNN & 421,738 & 93.03 & 92.69 $\pm$ 0.15 & 1.97 \\
LSTM / rows & 82,186 & 90.32 & 89.58 $\pm$ 0.05 & 1.77 \\
GRU / rows & 61,962 & 91.05 & 90.07 $\pm$ 0.17 & 1.73 \\
LSTM / patches & 76,042 & 88.83 & 87.65 $\pm$ 0.38 & 2.03 \\
GRU / patches & 57,354 & 88.97 & 88.00 $\pm$ 0.23 & 1.96 \\
\bottomrule
\end{tabularx}

\caption{E3 và hai baseline từ E1, trên cùng split.}
\end{table}
\tablenote

CNN khai thác lân cận hai chiều và chia sẻ kernel; LSTM/GRU tích lũy thông tin theo đường quét. Khi chuyển từ cuối một hàng patch sang đầu hàng kế tiếp, hai phần tử kề nhau trong chuỗi chưa chắc là hàng xóm gần nhất trên ảnh. Đây là đặc điểm của biểu diễn, không phải lỗi của cell.

Ảnh hưởng của cách quét được khảo sát bằng cách đánh giá checkpoint trên cùng ảnh validation với thứ tự gốc, đảo ngược và một hoán vị cố định. Mô hình giữ nguyên trọng số trong ba phép đo. Gradient theo vị trí chuỗi được đo đồng thời để mô tả đường truyền loss.
\IfFileExists{generated/e3_analysis.tex}{LSTM / rows trừ CNN E1: -3.11 $\pm$ 0.18 điểm phần trăm.

GRU / rows trừ CNN E1: -2.62 $\pm$ 0.13 điểm phần trăm.

LSTM / patches trừ CNN E1: -5.04 $\pm$ 0.35 điểm phần trăm.

GRU / patches trừ CNN E1: -4.69 $\pm$ 0.32 điểm phần trăm.

Rows và patches thay cả số bước lẫn số pixel/token; kết quả không tách riêng tác động độ dài chuỗi.
LSTM rows trừ patches: +1.93 $\pm$ 0.42 điểm phần trăm.

GRU rows trừ patches: +2.07 $\pm$ 0.22 điểm phần trăm.

GRU rows trừ LSTM rows: +0.49 $\pm$ 0.21 điểm phần trăm. Rows giữ trọn ngữ cảnh ngang trong mỗi token và có ít bước hơn patches; đây là một cách diễn giải phù hợp với kết quả trên cách chia dữ liệu này. Kiến trúc cell và đường quét cùng ảnh hưởng đến khả năng tổng hợp đặc trưng.
}{}

\resultfigure{e3_comparison_ce}{Validation CE của các mạng hồi quy và baseline E1 ở cùng seed.}{0.9}
\resultfigure{e3_comparison_accuracy}{So sánh accuracy của E3 với MLP/CNN.}{0.9}

\chapter{Mở rộng E3: Điều kiện truyền gradient}
\section{Vấn đề và giả thuyết}
Phần này khảo sát ảnh hưởng của khởi tạo, regularization và cách tổng hợp trạng thái lên LSTM. Chuỗi hàng và patch có độ dài lần lượt 28 và 49 bước, phù hợp để kiểm tra các thay đổi này trên bài toán ảnh nhỏ. Phân tích bắt đầu từ tích Jacobian của mạng hồi quy, sau đó đối chiếu với gradient và kết quả thực nghiệm.

Với trạng thái hồi quy tổng quát $s_t=F_\theta(s_{t-1},x_t)$ và $J_t=\partial s_t/\partial s_{t-1}$, đóng góp từ loss cuối chuỗi thỏa
\begin{equation}
\frac{\partial\ell_T}{\partial s_t}
=J_{t+1}^{\top}J_{t+2}^{\top}\cdots J_T^{\top}
\frac{\partial\ell_T}{\partial s_T},\qquad
\norm{\frac{\partial\ell_T}{\partial s_t}}
\leq\prod_{k=t+1}^{T}\norm{J_k}_2\,
\norm{\frac{\partial\ell_T}{\partial s_T}}.
\end{equation}
Tích nhiều phép co có thể làm gradient nhỏ; tích nhiều phép khuếch đại có thể làm nó lớn. Cận theo norm là một cách phân tích, không mô tả chính xác mọi hướng gradient. Với LSTM, $s_t$ cần gồm cả $h_t,c_t$, không chỉ một biến hidden \cite{pascanu}.

\section{Khởi tạo và regularization trực giao}
Với từng gate, ma trận hồi quy vuông $W_g\in\R^{H\times H}$ được khởi tạo trực giao ở các cấu hình tương ứng. Nếu $W_g^\top W_g=I$ thì $\norm{W_gv}_2=\norm{v}_2$: riêng phép biến đổi tuyến tính đó không làm thay đổi norm. Điều này không loại bỏ tác động của sigmoid, tanh, các gate và tích Jacobian toàn cell.

Tập các ma trận trực giao là trường hợp vuông của Stiefel manifold. Thay vì dùng optimizer có ràng buộc trên đa tạp, code dùng một penalty mềm trong không gian tham số thông thường:
\begin{equation}
\mathcal{R}_{orth}=\frac{1}{GH}\sum_{g=1}^{G}\norm{W_g^\top W_g-I}_F^2,
\qquad \mathcal{L}=\mathcal{L}_{CE}+\lambda\mathcal{R}_{orth},
\end{equation}
trong đó $G=4$ cho LSTM, $H=128$ và $\lambda=10^{-3}$ ở các cấu hình có regularization. Penalty được tính riêng cho bốn ma trận gate $H\times H$, rồi lấy trung bình.

Đặt $A=W^\top W-I=A^\top$. Vi phân cho một gate là
\begin{align}
d\,\norm{A}_F^2
&=2\operatorname{tr}(A\,dA)
 =2\operatorname{tr}\!\left[A(dW^\top W+W^\top dW)\right]\\
&=4\operatorname{tr}\!\left[(WA)^\top dW\right],
\end{align}
do đó
\begin{equation}
\nabla_{W_g}\mathcal{R}_{orth}=\frac{4}{GH}W_g(W_g^\top W_g-I).
\end{equation}
Công thức được kiểm tra với autograd. Penalty tạo chi phí tính toán thêm và có thể làm giảm mức linh hoạt của nghiệm. AdamW vẫn có weight decay cố định trên mọi tham số, nên nó đồng thời kéo trọng số về zero; không được mô tả quá trình train như duy trì trực giao chính xác.

Đây là penalty mềm trong không gian Euclid, không dùng diffusion maps, đồ thị dữ liệu hoặc optimizer manifold.

\section{Bias cổng nhớ theo độ dài chuỗi}
Ở cấu hình chrono, đặt $b_f=\log(T-1)$, $b_i=-\log(T-1)$; bias các gate khác bằng zero. Chỉ thay \code{bias_ih}; giữ \code{bias_hh=0} để tổng bias đúng giá trị dự định. Biến thể này dùng \emph{một timescale cố định}, thay cho cách lấy timescale ngẫu nhiên trong công trình chrono initialization gốc \cite{chrono}.

Khi bỏ qua đóng góp từ input/hidden tại lúc khởi tạo,
\begin{equation}
f_t\approx\sigma(\log(T-1))=1-\frac1T,
\qquad i_t\approx\frac1T.
\end{equation}
Giữ các gate cố định để chỉ xét đường trực tiếp $c_{t-1}\to c_t$, ta có
\begin{equation}
\left.\frac{\partial c_T}{\partial c_t}\right|_{\text{đường trực tiếp}}
=\prod_{k=t+1}^{T}\operatorname{diag}(f_k).
\end{equation}
Nếu mỗi $f_k\approx1-1/T$ thì độ suy giảm trên khoảng $T$ bước gần $(1-1/T)^T\approx e^{-1}$, thay vì $2^{-T}$ khi $f_k\approx0,5$. Đây chỉ là xấp xỉ để giải thích khởi tạo. Đạo hàm tổng còn có các đường qua hidden và đạo hàm của gate; không thể đồng nhất nó với tích trên.

Cổng quên gần 1 đi kèm cổng nhập nhỏ có thể khiến mô hình cập nhật thông tin mới chậm; bias lớn cũng có thể đưa sigmoid vào vùng bão hòa. Các bias tiếp tục được cập nhật trong quá trình huấn luyện; chrono được đánh giá như một cấu hình riêng.

\section{Mean pooling và clipping}
Baseline dự đoán từ $h_T$. Biến thể mean pooling dùng $\bar h=T^{-1}\sum_t h_t$ trước head. Bỏ qua dropout trong biểu thức minh họa, nếu $z=W_o\bar h+b$ và $y$ là vector one-hot thì đường gradient trực tiếp tới mỗi trạng thái là
\begin{equation}
\left.\frac{\partial\ell}{\partial h_t}\right|_{\text{head}}
=\frac1T W_o^\top(p-y).
\end{equation}
Đạo hàm tổng còn gồm các đường hồi quy về sau. Lợi ích tiềm năng là một trạng thái sớm không chỉ phụ thuộc vào việc tín hiệu đi qua toàn bộ chuỗi để tới head. Đổi lại, trung bình có thể làm mờ vị trí và trộn vùng nền vào biểu diễn. Learned position không được thêm cho E3; mạng hồi quy tự phụ thuộc thứ tự quét.

Clipping global norm được dùng cho tất cả baseline lẫn cải tiến:
\begin{equation}
g\leftarrow g\min\!\left(1,\frac{\tau}{\norm{g}_2}\right),\qquad \tau=1.
\end{equation}
Clipping giới hạn norm của gradient trước bước cập nhật, giúp kiểm soát exploding gradient. Cơ chế này không khuếch đại các gradient đang suy giảm \cite{pascanu}.

\section{Ablation và cách kiểm tra}
Các cấu hình dưới đây dùng LSTM theo hàng, giữ $H=128$, cách chia dữ liệu và ngân sách huấn luyện. ``Orth init'' thay khởi tạo. ``Orth reg'' thêm penalty vào bản orth init. ``Combined'' dùng trực giao ban đầu, penalty, chrono và mean pooling. Hai bản ``no reg'' và ``no mean'' bỏ riêng một thành phần từ combined; các thành phần khác giữ nguyên.

\begin{table}[htbp]\centering\small
\begin{tabularx}{\textwidth}{@{}Xrrrr@{}}
\toprule
Mô hình & Tham số & Val (\%) & Test (\%) & Giây/epoch \\
\midrule
LSTM / rows & 82,186 & 90.32 & 89.58 $\pm$ 0.05 & 1.77 \\
LSTM + orthogonal init & 82,186 & 90.20 & 89.35 $\pm$ 0.17 & 1.77 \\
LSTM + orthogonal init/reg & 82,186 & 90.07 & 89.47 $\pm$ 0.24 & 3.05 \\
LSTM + chrono bias & 82,186 & 90.23 & 89.28 $\pm$ 0.19 & 1.78 \\
LSTM + mean pooling & 82,186 & 89.28 & 88.32 $\pm$ 0.09 & 1.76 \\
LSTM combined & 82,186 & 89.17 & 87.92 $\pm$ 0.29 & 2.97 \\
Combined without regularizer & 82,186 & 89.07 & 87.91 $\pm$ 0.32 & 1.79 \\
Combined without mean pooling & 82,186 & 90.17 & 89.30 $\pm$ 0.23 & 3.03 \\
\bottomrule
\end{tabularx}

\caption{Baseline và bảy cấu hình mở rộng. Tất cả có cùng số tham số trainable.}
\end{table}
\tablenote

Chênh lệch orth init--baseline khảo sát khởi tạo; orth reg--orth init khảo sát penalty. Chrono và mean pooling được so riêng với baseline. Hai ablation no reg và no mean đánh giá đóng góp của từng thành phần trong tổ hợp. Chênh lệch trung bình và độ lệch chuẩn được tính từ các kết quả ghép theo cùng seed.

Phép đo unroll dùng trọng số của checkpoint native và công thức cell tương ứng, giữ gradient tại từng trạng thái trung gian. Loss dùng tổng theo batch để gradient từng mẫu không bị chia bởi batch size; sau đó lấy trung bình norm qua mẫu. Cách tính này cho phép khảo sát trực tiếp các trạng thái trong đồ thị BPTT.

Trên cùng 256 ảnh validation theo quy tắc định trước, diagnostic ghi $\norm{\partial\ell/\partial h_t}$, $\norm{\partial\ell/\partial c_t}$, gradient theo input token, giá trị gate trung bình và singular values của từng $W_g$. Các norm mô tả đường truyền loss của checkpoint đã học. Gradient nhỏ ở pixel nền cần được đọc cùng accuracy, CE và vị trí không gian của ảnh.
\IfFileExists{generated/improvement_analysis.tex}{Khởi tạo trực giao so với baseline làm thay đổi accuracy -0.23 $\pm$ 0.13 điểm phần trăm. Thêm penalty vào bản khởi tạo trực giao thay đổi +0.12 $\pm$ 0.13 điểm phần trăm, trong khi thời gian/epoch tăng từ 1.77s lên 3.05s.

LSTM + chrono bias so với baseline: -0.29 $\pm$ 0.19 điểm phần trăm.

LSTM + mean pooling so với baseline: -1.26 $\pm$ 0.11 điểm phần trăm.

Tổ hợp trừ baseline: -1.66 $\pm$ 0.30 điểm phần trăm. Bỏ mean pooling khỏi tổ hợp thay đổi accuracy +1.39 $\pm$ 0.52 điểm phần trăm; bỏ regularizer thay đổi -0.01 $\pm$ 0.35 điểm phần trăm và giảm thời gian/epoch 39.8\% theo các median đã đo. Mean pooling có ảnh hưởng rõ hơn trong hai ablation này; phần tính penalty chiếm thêm thời gian dù chênh lệch accuracy khi bỏ penalty khá nhỏ.

Cả bảy biến thể mở rộng có test accuracy trung bình thấp hơn baseline. Trong thiết lập này, các thay đổi chưa mang lại cải thiện accuracy. Mean pooling gộp đồng đều cả những trạng thái đầu chuỗi có ít ngữ cảnh; việc trộn các trạng thái này có thể làm loãng biểu diễn trước head.

LSTM / rows trên 256 ảnh validation: original accuracy 89.71\%; reversed accuracy 20.96\%; permuted accuracy 44.53\%. Hidden gradient norm trung bình tại bước 1/bước 28 là 0.293/0.304.

LSTM combined trên 256 ảnh validation: original accuracy 88.80\%; reversed accuracy 30.73\%; permuted accuracy 42.97\%. Hidden gradient norm trung bình tại bước 1/bước 28 là 0.34/0.0192.

Accuracy giảm khi đảo hoặc tráo thứ tự cho thấy checkpoint nhạy với thứ tự không gian. Mean pooling tạo đường gradient trực tiếp từ head tới mọi trạng thái hidden, trong khi last pooling nối head với trạng thái cuối; khác biệt này giải thích một phần hình dạng đường gradient của hai mô hình.
}{}

\resultfigure{improvement_comparison_accuracy}{Validation accuracy của các thay đổi chính ở cùng seed.}{0.9}
\resultfigure{e3_lstm_rows_gradients}{Gradient theo vị trí của LSTM baseline trên validation.}{0.88}
\resultfigure{e3_lstm_rows_combined_gradients}{Gradient của tổ hợp cải tiến, cùng cách chọn validation subset.}{0.88}

\chapter{Tái lập và giới hạn}
\section{Mã nguồn và tệp kết quả}
Repo tách mô hình thành \code{src/E1/}, \code{src/E2/}, \code{src/E3/}; phần dùng chung nằm trong \code{src/common/}. Mỗi run có \code{run.json}, \code{history.csv}, \code{best.pt}, \code{last.pt} và \code{metrics.json}. Đánh giá test sinh thêm prediction CSV, confusion matrix, ảnh sai và chỉ số theo lớp. Hyperparameter được khai báo trong một file JSON gồm 20 preset.

Checkpoint cuối epoch lưu model, AdamW, cosine scheduler, lịch sử và trạng thái RNG của Python, NumPy, PyTorch, CUDA, shuffle generator. Resume từ epoch hoàn tất gần nhất, không từ batch đang dở. Fingerprint kiểm tra code, config, seed, dữ liệu và normalization; môi trường torch/CUDA cũng được đối chiếu. Khóa thư mục ngăn hai tiến trình ghi cùng run. Cơ chế này không hứa khả năng tái lập bitwise giữa các phần cứng khác nhau.

\code{scripts/export_report.py} tổng hợp kết quả thành bảng LaTeX, CSV/JSON và hình vẽ. Chương trình đối chiếu cấu hình, seed, split và checkpoint trước khi gộp các lượt chạy.

\section{Cấu hình và các lượt chạy}
Thực nghiệm hoàn tất 60 lượt huấn luyện và 60 lượt đánh giá test. Bảng dưới liệt kê các seed của từng cấu hình.

{\small\begin{longtable}{@{}p{0.65\textwidth}rr@{}}
\toprule
Cấu hình & Seed train & Seed test \\ \midrule \endhead
Softmax & 42,43,44 & 42,43,44 \\
MLP & 42,43,44 & 42,43,44 \\
CNN & 42,43,44 & 42,43,44 \\
patches / manual & 42,43,44 & 42,43,44 \\
patches / torch & 42,43,44 & 42,43,44 \\
rows / manual & 42,43,44 & 42,43,44 \\
rows / torch & 42,43,44 & 42,43,44 \\
cnn / manual & 42,43,44 & 42,43,44 \\
cnn / torch & 42,43,44 & 42,43,44 \\
LSTM / rows & 42,43,44 & 42,43,44 \\
GRU / rows & 42,43,44 & 42,43,44 \\
LSTM / patches & 42,43,44 & 42,43,44 \\
GRU / patches & 42,43,44 & 42,43,44 \\
LSTM + orthogonal init & 42,43,44 & 42,43,44 \\
LSTM + orthogonal init/reg & 42,43,44 & 42,43,44 \\
LSTM + chrono bias & 42,43,44 & 42,43,44 \\
LSTM + mean pooling & 42,43,44 & 42,43,44 \\
LSTM combined & 42,43,44 & 42,43,44 \\
Combined without regularizer & 42,43,44 & 42,43,44 \\
Combined without mean pooling & 42,43,44 & 42,43,44 \\
\bottomrule
\end{longtable}
}

Các kiểm thử bao gồm forward/gradient của MSA, thứ tự token, unroll LSTM/GRU, đạo hàm regularizer, bias chrono, split và khôi phục trạng thái khi resume. Môi trường chạy, log và checksum được ghi trong \code{docs/TRAINING_EXECUTION.md} cùng thư mục kết quả.

Ba seed cho phép quan sát biến thiên giữa các lượt khởi tạo. Phạm vi so sánh giới hạn ở một tập dữ liệu, ngân sách 30 epoch và các siêu tham số đã chọn. Các chuỗi 28/49 bước cũng ngắn hơn nhiều so với những bài toán phụ thuộc dài hạn.

\section{Kết luận}
Ba nhóm mô hình được hiện thực và đánh giá trên cùng Fashion-MNIST. Các đối chứng cho thấy ảnh hưởng của biểu diễn đầu vào, cơ chế xử lý và điều kiện khởi tạo trong cùng ngân sách huấn luyện.
\IfFileExists{generated/conclusion.tex}{Trong 13 cấu hình cơ bản, CNN có test accuracy trung bình cao nhất, đạt 92.69\%. CNN-stem cho kết quả tốt nhất trong nhóm E2, còn biểu diễn theo hàng phù hợp hơn patch đối với cả LSTM và GRU. Bảy biến thể mở rộng E3 đều có accuracy trung bình dưới baseline LSTM rows. Các ablation cho thấy mean pooling ảnh hưởng lớn đến accuracy của tổ hợp, trong khi penalty trực giao làm tăng chi phí huấn luyện.
}{}

Báo cáo sử dụng AI được đính kèm tại tệp \code{AI_USE_REPORT.md} trong repository.

\begin{thebibliography}{9}
\addcontentsline{toc}{chapter}{Tài liệu tham khảo}
\bibitem{assignment} Lê Thành Sách, Nguyễn Quốc Minh. \emph{Bài tập E1–E3: Học sâu và ứng dụng trong thị giác máy tính, CO5085}. Đề học kỳ 261 (2026–2027).
\bibitem{fashion} H. Xiao, K. Rasul, R. Vollgraf. \emph{Fashion-MNIST: a Novel Image Dataset for Benchmarking Machine Learning Algorithms}. arXiv:1708.07747, 2017. \url{https://arxiv.org/abs/1708.07747}. Dữ liệu: \url{https://github.com/zalandoresearch/fashion-mnist}.
\bibitem{attention} A. Vaswani và cộng sự. \emph{Attention Is All You Need}. NeurIPS, 2017. \url{https://arxiv.org/abs/1706.03762}.
\bibitem{pascanu} R. Pascanu, T. Mikolov, Y. Bengio. \emph{On the difficulty of training recurrent neural networks}. ICML, 2013; arXiv:1211.5063. \url{https://arxiv.org/abs/1211.5063}.
\bibitem{chrono} C. Tallec, Y. Ollivier. \emph{Can recurrent neural networks warp time?} ICLR, 2018. \url{https://arxiv.org/abs/1804.11188}.
\bibitem{pytorch-lstm} PyTorch. \emph{LSTM: API and equations}. \url{https://docs.pytorch.org/docs/2.9/generated/torch.nn.LSTM.html}.
\bibitem{pytorch-gru} PyTorch. \emph{GRU: API and implementation notes}. \url{https://docs.pytorch.org/docs/2.9/generated/torch.nn.GRU.html}.
\bibitem{pytorch-versions} PyTorch. \emph{Previous versions: torch 2.10.0 / torchvision 0.25.0}. \url{https://pytorch.org/get-started/previous-versions/}.
\end{thebibliography}
\end{document}
